\chapter{Branches: merging and rebasing}
\label{ch:branches}

\chaptergoals{2 of 6 in Part~I.}
{\cref{ch:git}: your \file{git-playground} repository with a few commits.}
{develop on a branch, bring two lines of work together by merging or by
rebasing, know which of the two SHiP uses and why, and resolve a conflict.}

\section{Branches}

A \emph{branch} is just a movable name for a commit. When you commit on a
branch, the name moves forward to the new commit. \texttt{HEAD} is the branch
you are on. Work on anything that is not finished happens on a branch of its
own, so that \texttt{main} always works:

\begin{shell}
git switch -c feature          # create the branch "feature" and switch to it
# ... edit, add, commit ...
git switch main                # back to main; your files change accordingly
git branch                     # list branches; * marks the current one
\end{shell}

\begin{exercise}{A branch of your own}
In \file{git-playground}, create a branch \texttt{feature}, add a file
\file{feature.py} and commit it. Switch back to \texttt{main}: where did
\file{feature.py} go? Switch to \texttt{feature} again. Then look at both
branches at once with \cmd{git log -{}-oneline -{}-graph -{}-all}.
\end{exercise}

Meanwhile other work lands on \texttt{main}, and the two branches
\emph{diverge}. Sooner or later your branch has to be combined with
\texttt{main}. There are two ways to do that.

\section{Merging and rebasing}
\label{sec:mergerebase}

\begin{figure}[ht]
\centering
\begin{tikzpicture}[node distance=3.5mm]
  \newcommand{\panel}[1]{\node[font=\sffamily\small\bfseries, anchor=west] at (-0.35,1.25) {#1};}
  % ---- before ----
  \begin{scope}
    \panel{Before}
    \node[commit] (a) at (0,0) {A};
    \node[commit, right=of a] (b) {B};
    \node[commit, right=of b] (e) {E};
    \node[commit] (c) at ($(b)+(0.5,-1)$) {C};
    \node[commit, right=of c] (d) {D};
    \draw[parent] (b) -- (a); \draw[parent] (e) -- (b);
    \draw[parent] (c) -- (b); \draw[parent] (d) -- (c);
    \node[ref, above=1.2mm of e] {main};
    \node[ref, below=1.2mm of d] {feature};
  \end{scope}
  % ---- merge ----
  \begin{scope}[xshift=46mm]
    \panel{git merge feature}
    \node[commit] (a) at (0,0) {A};
    \node[commit, right=of a] (b) {B};
    \node[commit, right=of b] (e) {E};
    \node[commit] (c) at ($(b)+(0.5,-1)$) {C};
    \node[commit, right=of c] (d) {D};
    \node[mergecommit, right=of e] (mm) {M};
    \draw[parent] (b) -- (a); \draw[parent] (e) -- (b);
    \draw[parent] (c) -- (b); \draw[parent] (d) -- (c);
    \draw[parent] (mm) -- (e); \draw[parent] (mm) -- (d);
    \node[ref, above=1.2mm of mm] {main};
    \node[ref, below=1.2mm of d] {feature};
  \end{scope}
  % ---- rebase ----
  \begin{scope}[xshift=100mm]
    \panel{git rebase main}
    \node[commit] (a) at (0,0) {A};
    \node[commit, right=of a] (b) {B};
    \node[commit, right=of b] (e) {E};
    \node[rewritten, right=of e] (c2) {C$'$};
    \node[rewritten, right=of c2] (d2) {D$'$};
    \draw[parent] (b) -- (a); \draw[parent] (e) -- (b);
    \draw[parent] (c2) -- (e); \draw[parent] (d2) -- (c2);
    \node[ref, above=1.2mm of e] {main};
    \node[ref, above=1.2mm of d2] {feature};
  \end{scope}
\end{tikzpicture}
\caption{Combining a branch with \texttt{main}. Arrows point from each commit
to its parent. A merge adds a commit M with two parents. A rebase replays C and
D on top of E as new commits C$'$ and D$'$ (dashed): same changes, new hashes,
and a straight history; \texttt{main} can then simply move forward to D$'$
(a \emph{fast-forward}).}
\label{fig:mergerebase}
\end{figure}

\paragraph{Merging} (\cmd{git merge feature}, run on \texttt{main}) keeps
both lines of history exactly as they happened and joins them with a new
\emph{merge commit} that has two parents. Nothing is rewritten. The price is a
history that branches and rejoins everywhere, which is hard to read and to
bisect.

\paragraph{Rebasing} (\cmd{git rebase main}, run on \texttt{feature}) takes
the commits that are on \texttt{feature} but not on \texttt{main}, and
re-applies them one by one on top of the current \texttt{main}. The result
looks as if you had started your work from today's \texttt{main}: a straight
line, with no merge commits. The price is that your commits are
\emph{replaced} by new ones with new hashes.

\begin{center}\small
\begin{tabularx}{\linewidth}{l L L}
\toprule
 & merge & rebase \\
\midrule
history     & as it happened, with merge commits & straight line \\
your commits& kept, same hashes & rewritten, new hashes \\
conflicts   & resolved once, in the merge commit & resolved per commit, as each is replayed \\
safe on shared branches & yes & \textbf{no}: others still have the old commits \\
\bottomrule
\end{tabularx}
\end{center}

\paragraph{What SHiP uses.} The \texttt{main} branch of the ShipSoft
repositories requires a \emph{linear history}: merge commits are not
accepted. Pull requests land by rebasing (or by squashing all their commits
into one), through a merge queue (\cref{sec:pullrequests}). So when your
branch falls behind \texttt{main}, you bring it up to date by \emph{rebasing
it}, not by merging \texttt{main} into it.

\begin{pitfall}
The golden rule: only rebase commits that nobody else has built on. Your own
pull-request branch is fine; \texttt{main}, or a branch a colleague is also
working on, is not. After rebasing a branch you had already pushed, push with
\cmd{git push -{}-force-with-lease} (\cref{ch:remotes}).
\end{pitfall}

\begin{exercise}{Merge, then rebase}
Continue from the previous exercise: switch to \texttt{main} and make one
commit there, so the branches diverge. Copy the branch
(\cmd{git branch feature-copy feature}). Merge \texttt{feature} into
\texttt{main} and look at \cmd{git log -{}-oneline -{}-graph -{}-all}. Undo the
merge with \cmd{git reset -{}-hard HEAD\textasciitilde1} (only ever on a
playground), switch to \texttt{feature-copy}, run \cmd{git rebase main}, and
look at the graph again. Compare the hashes of the rebased commits with the
originals. Finally, \cmd{git switch main} and
\cmd{git merge -{}-ff-only feature-copy}: why does this one not create a merge
commit?
\end{exercise}

\section{Conflicts}

When both sides changed the same lines, git cannot decide and stops. During a
rebase it looks like this:

\begin{outputlisting}
CONFLICT (content): Merge conflict in mystery_analysis.py
error: could not apply c8dfe70... feat(task-10): log-scale plot
hint: Resolve all conflicts manually, mark them as resolved with
hint: "git add/rm <conflicted_files>", then run "git rebase --continue".
\end{outputlisting}

and the file contains both versions between markers:

\begin{outputlisting}
<<<<<<< HEAD
OUTPUT_PDF = "results/mystery_fit.pdf"
=======
OUTPUT_PDF = "mystery_fit_log.pdf"
>>>>>>> c8dfe70 (feat(task-10): log-scale plot)
\end{outputlisting}

Edit the file into what it should be, delete the three marker lines, then
\begin{shell}
git add mystery_analysis.py
git rebase --continue          # or: git rebase --abort, to go back to before
\end{shell}

\begin{pitfall}
During a rebase the sides are the other way round from what you might expect:
\texttt{HEAD} (the upper half) is the branch you are rebasing \emph{onto},
i.e.\ the current \texttt{main}; the lower half is your own commit being
replayed. Always run the tests again after resolving a conflict: a resolution
that git accepts is not necessarily one that works.
\end{pitfall}

\begin{exercise}{A conflict on purpose}
Two branches change the same line of the same file: \texttt{main} and
\texttt{mine} each give \file{settings.py} a different output name.
\begin{shell}
vim settings.py                # write: OUTPUT = "plot.pdf"
git add settings.py && git commit -m "feat: settings"
git switch -c mine
vim settings.py                # change to: OUTPUT = "mine.pdf"
git commit -am "feat: my output"
git switch main
vim settings.py                # change to: OUTPUT = "theirs.pdf"
git commit -am "feat: their output"
git switch mine
git rebase main                # stops: CONFLICT in settings.py
\end{shell}
Open \file{settings.py} in your editor, decide what the line should be, remove the markers,
\cmd{git add settings.py} and \cmd{git rebase -{}-continue} (git opens an
editor for the message: save and close it). Look at the graph. Then make the
same conflict again and leave it with \cmd{git rebase -{}-abort} instead.
\end{exercise}

\begin{levelcheck}
\begin{itemize}
  \can create, switch between and list branches;
  \can explain, with a drawing, what a merge and a rebase do to the history;
  \can say why ShipSoft wants rebased, linear history, and which branches you
       may and may not rebase;
  \can resolve a conflict during a rebase, or abort it.
\end{itemize}
\end{levelcheck}
